\section*{Chapter 2 --- Figure 1: Transport plan vs Brenier map}
 
% --- Figure 3: Transport plan vs Brenier map ---
\begin{figure}[h]
\centering
\begin{tikzpicture}[font=\small, scale=0.92]
% LEFT: Kantorovich
\begin{scope}[shift={(-3.4,0)}]
  \node[font=\small\bfseries, deepblue, align=center]
    at (0,3.4) {Kantorovich\\transport plan $\gamma$};
  \foreach \x/\h in {-1.6/0.55,-0.8/1.05,0/1.35,0.8/1.0,1.6/0.60}{
    \fill[deepblue!65,opacity=0.85]
      (\x-0.26,-2.5) rectangle (\x+0.26,-2.5+\h);
    \draw[deepblue!80](\x-0.26,-2.5)rectangle(\x+0.26,-2.5+\h);
  }
  \node[deepblue,font=\footnotesize\bfseries] at (0,-2.92) {$\mu$};
  \foreach \x/\h in {-1.2/0.55,-0.4/0.95,0.4/1.25,1.2/0.85,1.9/0.45}{
    \fill[softgreen!60,opacity=0.85]
      (\x-0.26,2.5) rectangle (\x+0.26,2.5-\h);
    \draw[softgreen!70!black](\x-0.26,2.5)rectangle(\x+0.26,2.5-\h);
  }
  \node[softgreen!80!black,font=\footnotesize\bfseries]
    at (0.35,2.92) {$\nu$};
  \foreach \xs/\xt in {
    -1.6/-1.2,-1.6/-0.4,-1.6/0.4,
    -0.8/-1.2,-0.8/-0.4,-0.8/0.4,
     0.0/-0.4, 0.0/0.4, 0.0/1.2,
     0.8/-0.4, 0.8/0.4, 0.8/1.2,
     1.6/0.4,  1.6/1.2, 1.6/1.9
  }{\draw[gray!35,thin](\xs,-1.5)--(\xt,1.5);}
  \node[gray!55,font=\scriptsize,align=center] at (0.1,0.05)
    {$\gamma(x,y)$\\mass splitting\\allowed};
\end{scope}
\draw[gray!35,thick] (0,-3.2)--(0,3.2);
% RIGHT: Brenier
\begin{scope}[shift={(3.4,0)}]
  \node[font=\small\bfseries, deepblue, align=center]
    at (0,3.4) {Brenier\\map $T=\nabla\varphi$};
  \foreach \x/\h in {-1.6/0.55,-0.8/1.05,0/1.35,0.8/1.0,1.6/0.60}{
    \fill[deepblue!65,opacity=0.85]
      (\x-0.26,-2.5) rectangle (\x+0.26,-2.5+\h);
    \draw[deepblue!80](\x-0.26,-2.5)rectangle(\x+0.26,-2.5+\h);
  }
  \node[deepblue,font=\footnotesize\bfseries] at (0,-2.92) {$\mu$};
  \foreach \x/\h in {-1.2/0.55,-0.4/0.95,0.4/1.25,1.2/0.85,1.9/0.45}{
    \fill[softgreen!60,opacity=0.85]
      (\x-0.26,2.5) rectangle (\x+0.26,2.5-\h);
    \draw[softgreen!70!black](\x-0.26,2.5)rectangle(\x+0.26,2.5-\h);
  }
  \node[softgreen!80!black,font=\footnotesize\bfseries]
    at (0.35,2.92) {$\nu$};
  \foreach \xs/\xt in {-1.6/-1.2,-0.8/-0.4,0.0/0.4,0.8/1.2,1.6/1.9}{
    \draw[->,deepblue!55, dashed, line width=1.1pt]
      (\xs,-1.45)--(\xt,1.45);
  }
  \node[deepblue!70,font=\scriptsize,align=center] at (0.15,0.05)
    {$T(x)=\nabla\varphi(x)$\\deterministic,\\no mass splitting};
\end{scope}
\end{tikzpicture}
\caption{Transport plan (Kantorovich) versus deterministic map
  (Brenier). Kantorovich allows mass splitting (faint lines,
  left). Under Brenier's assumptions, the $W_2$-optimal
  transport is realized by a deterministic gradient map
  $T=\nabla\varphi$: no mass splitting, no crossing (right).}
\end{figure}
 
\clearpage
\section*{Chapter 2 --- Figure 2: Displacement interpolation}
 
% --- Figure 4: Displacement interpolation / space-time diagram ---
\begin{figure}[h]
\centering
\begin{tikzpicture}[font=\small, scale=0.95]
 
% ---- Top row: density snapshots ----
\foreach \t/\xoff in {0/0, 0.5/4.3, 1/8.6}{
  \begin{scope}[shift={(\xoff,5.3)}]
    \draw[gray!30] (-1.7,-0.1) rectangle (1.7,1.5);
    \pgfmathsetmacro{\mlo}{-1.0+1.5*\t}
    \pgfmathsetmacro{\mhi}{0.6+1.0*\t}
    \fill[deepblue!55,opacity=0.85,domain=-1.6:1.6,samples=50,
          variable=\x]
      plot(\x,{0.05+1.3*exp(-3.0*(\x-(\mlo+0.3*(\mhi-\mlo)))*
                (\x-(\mlo+0.3*(\mhi-\mlo))))})
      -- (1.6,0.05) -- (-1.6,0.05) -- cycle;
    \node[gray!60,font=\scriptsize] at (0,-0.4) {$t=\t$};
  \end{scope}
}
\node[deepblue,font=\small\bfseries] at (4.3,6.1)
  {Mass moves, it does not fade};
 
% ---- Bottom: space-time diagram ----
\begin{scope}[shift={(0,0)}]
  \draw[->,thick,gray!50] (0,0) -- (9.0,0)
    node[right,font=\scriptsize]{space};
  \draw[->,thick,gray!50] (0,0) -- (0,4.4)
    node[above,font=\scriptsize]{time};
 
  \foreach \xs/\xt in {
    1.0/3.5, 1.6/4.6, 2.2/5.7, 2.8/6.5, 3.4/7.3
  }{
    \draw[deepblue!55,thick] (\xs,0) -- (\xt,4.0);
  }
 
  \filldraw[deepblue] (1.0,0) circle (2pt);
  \filldraw[deepblue] (1.6,0) circle (2pt);
  \filldraw[deepblue] (2.2,0) circle (2pt);
  \filldraw[deepblue] (2.8,0) circle (2pt);
  \filldraw[deepblue] (3.4,0) circle (2pt);
 
  \filldraw[softgreen!70!black] (3.5,4.0) circle (2pt);
  \filldraw[softgreen!70!black] (4.6,4.0) circle (2pt);
  \filldraw[softgreen!70!black] (5.7,4.0) circle (2pt);
  \filldraw[softgreen!70!black] (6.5,4.0) circle (2pt);
  \filldraw[softgreen!70!black] (7.3,4.0) circle (2pt);
 
  \node[deepblue,font=\scriptsize] at (1.0,-0.35) {$\mu$};
  \node[softgreen!70!black,font=\scriptsize] at (7.3,4.35) {$\nu$};
 
  \draw[warmgold,dashed,thick] (0,2.0) -- (9.0,2.0)
    node[midway,above=2pt,font=\scriptsize,warmgold!80!black]
      {$\rho_{1/2}$: displacement interpolation};
 
  \node[deepblue!70, font=\scriptsize, align=center,
        draw=deepblue!25, rounded corners=2pt,
        fill=white, inner sep=4pt] at (8.0,1.0)
    {$W_2^2(\mu,\nu)$\\$=\int_0^1\!\int|v_t|^2\rho_t\,dx\,dt$};
\end{scope}
 
\end{tikzpicture}
\caption{Displacement interpolation: the cheapest movie. Top:
  density snapshots at $t=0, 0.5, 1$ show mass moving, not
  fading. Bottom: the space-time diagram of straight-line
  characteristic paths from $\mu$ to $\nu$, with the
  displacement interpolation $\rho_{1/2}$ marked. Chapter~4
  revisits this diagram with stochastic paths.}
\end{figure}
 
\end{document}
 
